\section{Simulación}

\subsection{Materiales}

\begin{itemize}
    \item Computadora o laptop
    \item Software GeoGebra
    \item Cuaderno de apuntes
    \item Calculadora científica
\end{itemize}

\subsection{Procedimiento}

\begin{enumerate}
    \item Abrir el software GeoGebra.
    \item En la barra de entrada, definir la función de posición del automóvil (MRU):
    \[
    x_a(t) = 40t
    \]
    \item Definir la función de posición del camión (MRUA):
    \[
    x_c(t) = 2t^2 + 80
    \]
    \item Observar las curvas generadas: la del automóvil es una línea recta y la del camión es una parábola.
    \item Identificar gráficamente los puntos de intersección entre ambas curvas.
    \item Verificar los resultados obtenidos analíticamente con los valores leídos en la gráfica.
\end{enumerate}

El ejercicio simulado es el siguiente:

\textbf{Un automóvil se mueve a velocidad constante de 144 km/h en una carretera recta. En determinado momento, un camión que está a 80 m delante del automóvil comienza a moverse desde el reposo con una aceleración de 4 m/s$^2$. Determinar:}
\begin{itemize}
    \item[{a)}] El tiempo y la posición cuando el automóvil y el camión se encuentran juntos.
    \item[{b)}] El momento cuando la velocidad del automóvil y del camión son iguales, y la separación entre ellos en ese instante.
\end{itemize}

\subsection{Resolución}

\textbf{Datos del problema:}

\begin{itemize}
    \item \textbf{Automóvil:}
    \begin{align*}
        v_a &= 144\,\frac{\cancel{\text{km}}}{\cancel{\text{h}}} \cdot \frac{1000\,\text{m}}{1\,\cancel{\text{km}}} \cdot \frac{1\,\cancel{\text{h}}}{3600\,\text{s}} = 40\,\frac{\text{m}}{\text{s}} \\
        x_{0a} &= 0\,\text{m} \\
        a_a &= 0\,\frac{\text{m}}{\text{s}^2}
    \end{align*}

    \item \textbf{Camión:}
    \begin{align*}
        v_{0c} &= 0\,\frac{\text{m}}{\text{s}} \\
        x_{0c} &= 80\,\text{m} \\
        a_c &= 4\,\frac{\text{m}}{\text{s}^2}
    \end{align*}
\end{itemize}

\textbf{Velocidad del camión (MRUA):}
\vspace{-20pt}

\begin{align*}
a_c &= \frac{dv}{dt} \\
4 &= \frac{dv}{dt} \\
dv &= 4\,dt \\
\int_{0}^{v} dv &= \int_{0}^{t} 4\,dt \\
\left[v\right]_{0}^{v} &= 4\left[t\right]_{0}^{t} \\
v - 0 &= 4(t - 0) \\
\Aboxed{v_c &= 4t}
\end{align*}

\textbf{Posición del camión (MRUA):}
\vspace{-20pt}

\begin{align*}
v_c &= \frac{dx}{dt} \\
4t &= \frac{dx}{dt} \\
dx &= 4t\,dt \\
\int_{80}^{x} dx &= \int_{0}^{t} 4t\,dt \\
\left[x\right]_{80}^{x} &= 2\left[t^2\right]_{0}^{t} \\
x - 80 &= 2(t^2 - 0) \\
x - 80 &= 2t^2 \\
\Aboxed{x_c &= 2t^2 + 80}
\end{align*}

\textbf{Posición del automóvil (MRU):}
\vspace{-20pt}

\begin{align*}
v_a &= \frac{dx}{dt} \\
40 &= \frac{dx}{dt} \\
dx &= 40\,dt \\
\int_{0}^{x} dx &= \int_{0}^{t} 40\,dt \\
\left[x\right]_{0}^{x} &= 40\left[t\right]_{0}^{t} \\
x - 0 &= 40(t - 0) \\
\Aboxed{x_a &= 40t}
\end{align*}

\subsubsection*{a) Tiempo y posición cuando se encuentran juntos}

Para que el automóvil y el camión estén en la misma posición, igualamos sus ecuaciones de posición:
\vspace{-20pt}

\begin{align*}
x_a &= x_c \\
40t &= 2t^2 + 80 \\
0 &= 2t^2 - 40t + 80 \\
0 &= t^2 - 20t + 40
\end{align*}

Aplicando la fórmula general:
\vspace{-20pt}

\begin{align*}
t &= \frac{-b \pm \sqrt{b^2 - 4ac}}{2a} \\
t &= \frac{20 \pm \sqrt{(-20)^2 - 4(1)(40)}}{2(1)} \\
t &= \frac{20 \pm \sqrt{400 - 160}}{2} \\
t &= \frac{20 \pm \sqrt{240}}{2} \\
t &= \frac{20 \pm 15{,}49}{2}
\end{align*}

Se obtienen dos eventos:
\vspace{-20pt}

\[
\textbf{Primer evento:} \quad t_1 = \frac{20 - 15{,}49}{2} = \frac{4{,}51}{2} \approx 2{,}25\,\text{s} \Rightarrow t_1 \approx 2{,}3\,\text{s}
\]

\[
\textbf{Segundo evento:} \quad t_2 = \frac{20 + 15{,}49}{2} = \frac{35{,}49}{2} \approx 17{,}74\,\text{s}
\]

\textit{Nota: El primer evento ($t_1 \approx 2{,}3$ s) es físicamente el relevante; corresponde al momento en que el auto alcanza al camión por primera vez.}

Calculando la posición en $t_1 \approx 2{,}3$ s:

\begin{align*}
x_a &= 40t_1 = 40(2{,}25) = 90\,\text{m} \\
x_c &= 2(2{,}25)^2 + 80 = 2(5{,}0625) + 80 = 10{,}125 + 80 \approx 90{,}2\,\text{m}
\end{align*}

\[
\boxed{t_1 \approx 2{,}3\,\text{s}, \qquad x \approx 90{,}2\,\text{m}}
\]

\subsubsection*{b) Momento en que las velocidades son iguales y separación}

\vspace{-40pt}

\begin{align*}
\intertext{Igualamos la velocidad del automóvil con la del camión:}
v_a &= v_c \\
40 &= 4t \\
t &= \frac{40}{4} \\
\Aboxed{t &= 10\,\text{s}}
\intertext{Evaluamos las posiciones en $t = 10$ s:}
x_a &= 40t = 40(10) = 400\,\text{m} \\
x_c &= 2t^2 + 80 = 2(10)^2 + 80 = 200 + 80 = 280\,\text{m}
\intertext{Separación entre ellos:}
\Delta x &= x_a - x_c \\
\Delta x &= 400 - 280 \\
\Aboxed{\Delta x &= 120\,\text{m}}
\end{align*}

\subsection{Resultados}

A partir de la simulación realizada en GeoGebra se obtuvieron los siguientes resultados:

\begin{itemize}
    \item La gráfica del automóvil es una línea recta (MRU), con pendiente constante igual a su velocidad (40 m/s).
    \item La gráfica del camión es una parábola (MRUA), que parte desde $x = 80$ m y se acelera progresivamente.
    \item Las dos curvas se intersectan en dos puntos: el primero alrededor de $t \approx 2{,}3$ s y el segundo en $t \approx 17{,}74$ s.
\end{itemize}

\begin{table}[H]
\centering
\begin{tabular}{c|c|c|c}
Tiempo (s) & $x_{\text{auto}}$ (m) & $x_{\text{camión}}$ (m) & $v_{\text{camión}}$ (m/s) \\\hline
0   & 0     & 80    & 0  \\
2{,}3 & 90  & 90{,}2 & 9{,}2 \\
5   & 200   & 130   & 20 \\
10  & 400   & 280   & 40 \\
17{,}74 & 709{,}6 & 709{,}6 & 70{,}96 \\
\end{tabular}
\caption{Datos obtenidos en la simulación}
\end{table}

\vspace{-20pt}

\subsection{Análisis}

Los resultados obtenidos permiten observar que:

\begin{itemize}
    \item El automóvil, al tener mayor velocidad inicial, alcanza rápidamente al camión en $t \approx 2{,}3$ s (posición $\approx 90{,}2$ m). Sin embargo, como el camión está acelerando, eventualmente lo supera, por lo que el automóvil vuelve a alcanzarlo en el segundo evento ($t \approx 17{,}74$ s).
    \item En $t = 10$ s las velocidades son iguales (ambas 40 m/s), pero el camión aún está 120 m por detrás del automóvil. A partir de ese momento, el camión, al seguir acelerando, comenzará a reducir esa brecha.
    \item La diferencia en el tipo de movimiento (MRU vs MRUA) genera dos intersecciones en la gráfica posición-tiempo, lo cual tiene un significado físico claro: primero el auto pasa al camión, y luego el camión recupera y alcanza al auto.
\end{itemize}